\part{Fast Writing}

\chapter{The Measured Hand}

Before writing speed could be trained it had to be measured, and before it could be measured it needed a unit. The history of typing tutors begins with the history of that unit.

\section{Copying as Labor}

For most of the history of the codex, reproducing a text meant copying it by hand, and copying was paid work. Payment by the piece is one of the oldest ways of turning a bodily rate into an economic quantity. The university stationers of medieval Europe are a familiar example: exemplars were divided into sections (the \emph{pecia} system) and rented out so that several scribes could copy one work at once. \needs{secondary source on the pecia system and its dates}

The point for this book is structural. Piece rates and rented exemplars already contain the two components of every later tutor: a model to be copied, and a rule for judging the copy. What they lack is an explicit measure of rate per unit time, since the scribe was paid for the product and not for the pace.

\section{The Unit of Speed}

A rate needs a countable unit. Natural words vary too much in length to serve, so later practice fixed a conventional word of five characters, spaces included. Under this convention gross speed is
\[
\wpm_{\text{gross}}=\frac{c/5}{t},
\]
where $c$ is the number of characters typed and $t$ the time in minutes. The convention matters because it converts a linguistic quantity, the word, into a mechanical one, the keystroke. Every timed test in this book, including the five-minute writings of Chapter~8, rests on it. Appendix~B states the scoring formulas and the variations found in the sources.

\section{Speed as Capacity}

Treating speed as a stable property of a person, something a test can reveal and practice can raise, is a modern idea. It requires standardized text, a clock, and an examiner. Each of these reappears in the chapters that follow: the standardized text as the timed passage, the clock as the metronome, the phonograph record and the cassette, and the examiner as the supervisor who initials a test sheet.

\chapter{Shorthand Traditions}

Shorthand is the oldest sustained attempt to design inscription around the limits of the moving hand. Its central move is to separate the recording of language from ordinary spelling.

\section{Ancient and Early Modern Systems}

The Tironian notes, associated with Marcus Tullius Tiro, the secretary of Cicero, are the classical reference point for abbreviated writing. \needs{dating and scope of the Tironian notes as a system versus a repertoire of signs} In England, Timothy Bright published a system of ``characterie'' in 1588, and later seventeenth-century systems, among them Thomas Shelton's, were in wide use; Shelton's system is the one associated with the diary of Samuel Pepys. \needs{verify Bright 1588 and the Shelton attribution}

\section{Phonetic Systems}

The nineteenth century brought systems built on sound. Isaac Pitman's phonography (1837) writes what is heard, with shading to distinguish voiced from unvoiced consonants. John Robert Gregg's light-line system (1888) reduced the dependence on shading and made outlines more cursive. \needs{confirm first-publication dates against a standard history}

\section{What Shorthand Teaches a Typing Tutor}

Shorthand instruction anticipated three features of later tutors.

\begin{itemize}
\item \textbf{Graded units.} Learners meet a small set of signs first and build words from them, so the lesson order is itself a specification of a growing symbol set.
\item \textbf{Dictation as pacing.} Speed is developed by reading aloud at a stated rate, which is the ancestor of the metronome, the phonograph record, and the cassette tape.
\item \textbf{The gap between transcript and record.} A shorthand note is not the text; it must be transcribed. That gap between a fast record and a finished document recurs whenever a system separates capture from rendering.
\end{itemize}

\chapter{Stenotype and Machine Shorthand}

Machine shorthand replaced the moving pencil with a keyboard on which several keys are pressed together. The machine records sound, not spelling.

\section{Development}

Machine shorthand is commonly traced to nineteenth-century American patents, with the twentieth-century Stenotype of Ward Stone Ireland (1911) as the design that established the pattern of a compact chorded keyboard printing on a paper strip. \needs{patent dates and inventors; distinguish Bartholomew, Ireland, and later commercial machines}

\section{Institutional Setting}

Court reporting is the paradigm application. The reporter's output is a paper strip in phonetic notation that must be transcribed into ordinary text, so the profession divides into capture and transcription, and its training divides accordingly into speed building and translation. \needs{training standards and certification speeds, by period}

\section{Scope of This Chapter}

This chapter is historical and institutional. The theoretical question, namely what happens to the notion of a keystroke when the primitive action is a chord, is reserved for Chapter~13, so that the two chapters divide the subject between them and do not repeat each other.

\chapter{Telegraphy}

Telegraph operators were the first large occupational group whose skill was defined by a sustained, measured rate of sending and receiving.

\section{Sending and Receiving}

Samuel Morse's telegraph demonstrated a working line in 1844. \needs{confirm date and route} Operators were tested on both sending and receiving. Receiving initially involved reading a paper record, and skilled operators came to receive by ear.

\section{A Labor Metric}

Telegraph speed tests fixed the practice of stating skill as a number of words per minute, with a conventional word length so that different messages could be compared. The convention traveled with the operators. Many early typewriter operators had telegraph backgrounds, and the typewriter was used to transcribe received messages directly. \needs{documentation of the telegrapher-to-typist pathway}

\section{Two Things Carried Forward}

Two features of telegraph training recur in later tutors. The first is training by pacing signal: the learner keeps up with a sender. The second is a distinction between copying at speed and copying without error, a distinction that Chapter~15 will show is written into the printed instructions of at least one typing workbook.

\chapter*{Study Questions, Part I}
\addcontentsline{toc}{chapter}{Study Questions, Part I}
\begin{enumerate}
\item Why does a rate of writing require a conventional unit of length? What does the five-character word gain and lose compared with counting natural words?
\item Shorthand and typing both separate capture from a finished document. Where does each place the cost of the gap?
\item Compare piece-rate payment for copying with a modern words-per-minute test. What is measured in each case, and who is the examiner?
\item Which features of telegraph training appear in the cassette-paced practice of Chapter~8?
\end{enumerate}
